\begin{frame}{Síntese e preparação para a próxima aula}
\begin{itemize}
\item O tipo deve representar a natureza e a faixa dos dados.
\item Declaração cria o objeto; inicialização estabelece seu primeiro valor.
\item A atribuição altera o valor e pode realizar conversões.
\item Nomes claros, unidades explícitas e formatos corretos tornam o programa verificável.
\item Tamanhos e limites precisam ser conferidos no ambiente.
\end{itemize}
\begin{block}{Bilhete de saída}
Explique a diferença entre \texttt{int nivel = 10;} e \texttt{nivel = 20;}. Cite um dado que você representaria por \texttt{bool}.
\end{block}
Próxima aula: operar, comparar e combinar os valores armazenados.
\end{frame}
